The memory of a computer contains nothing but zeros and ones. So we must be able to
represent all the information we want to store as sequences of zeros and
ones. This translation of information into binary is called {\em binary coding}. The digits
0 and 1 used to represent information are called {\em bits}. We shall
begin by studying other, older forms of coding.

\section{Morse code}
The electric telegraph, the ancestor of our telephone, was invented in 1837 by
the American \href{https://en.wikipedia.org/wiki/Samuel_Morse}{Samuel Morse} \cite{bertho}. The sending telegraph operator has a key, which is a simple spring
lever used to send short signals (called dots) and long signals (called
dashes) separated by silences or blanks (corresponding to the absence of a signal). The signal is
carried along an electric wire, the return path being through the earth. Each
movement of the key closes the circuit and causes a short or long movement of the receiver,
which reproduces the signals
on a paper tape. To send a message, one first shortens it into {\em telegraphic
style}, then translates each sign (letter, digit, punctuation mark, \ldots) into
a sequence of dots and dashes separated by blanks according to the {\em Morse code} of figure
\ref{fig:codemorse} (page \pageref{fig:codemorse}). Seven blanks are left at the end of each word.
Various other codes are used to mark the beginning and the end of the message, to signal errors,
etc. Morse's telegraph could transmit about 25 words a minute.
\begin{figure}[t]
\vspace{9.3cm}%
\includePostScript{CodeMorse}{0.95}{0.95}%
% \vspace{8.8cm} %8.9cm donne la figure sur une seule page !
 \caption{Morse code}\label{fig:codemorse}
\end{figure}

\vspace{5.0cm}

\section{Number systems}
\index{Decimal number system}
Our way of coding the natural numbers (positive or zero) is a positional, decimal number system.
It is a decimal system because we use the ten digits 0, 1, \ldots, 9, which are of
Indo-Arabic origin. It is a positional system because each digit in a number represents a
value that depends on its position in the number. For example in 1789, the 1 represents 1000,
the 7 represents 700, the 8 represents 80 and the 9 represents 9. So we have $1789 = 1000 + 700
+ 80 + 9 = (1 \times 10^{3}) +  (7 \times 10^{2}) +  (8 \times 10^{1}) + (9 \times 10^{0})$.
For the digits after the decimal point we use the inverses. For example $0.147= (1\times 10^{-1})+
(4\times 10^{-2})+(7\times 10^{-3})=\frac{1}{10}+\frac{4}{100}+\frac{7}{1000}=\frac{100}{1000}+\frac{40}{1000}+\frac{7}{1000}=\frac{147}{1000}$.
Mathematicians, ancient and modern, have invented many other ways of coding numbers.
Figure \ref{fig:code1-9} (page \pageref{fig:code1-9}) shows different ways of coding the number
123,456,789, some of them going back to ancient civilisations \cite{guitel}. The superiority of
the positional, decimal coding of the natural numbers does not come from the choice of base ten but
from the ease with which the usual arithmetic operations can be carried out,
for example with an abacus (long in use, and now gone from
primary schools).

\begin{figure}[htbp]
 \includevspacePostScript[18.0cm]{Codages123456789}{0.95}{0.95}
 \caption{Coding of 123,456,789 in Babylonian, Egyptian, Greek, Roman,\label{fig:code1-9}}
 \centering{hexadecimal, decimal, octal, ternary and binary number systems}
\end{figure}


\section{Binary coding}

All the information stored in the computer's memory or on disk must
be represented by sequences of zeros and ones. On microcomputers,
as they are designed today, the length of these sequences of bits must be a multiple
of 8. Depending on their length they are given special names, such as {\em byte} for a sequence
of 8 bits, {\em half-word} for a sequence of 16 bits, {\em word} for 32 bits and {\em double word}
for 64.

To show that any information can be represented by a sequence of bits (provided
we know the code), we shall take the example of integers, then that of characters and
of texts typed on the keyboard.

\subsection{Natural numbers}

\index{Binary number system}
To represent positive integers by a sequence of bits, it is enough to write them in base
2. For example:
\[\begin{array}{lll}
147 & = &128 +  16 + 2 + 1 \\
       & = & (1 \times 2^{7}) + (0 \times 2^{6}) + (0 \times 2^{5}) + (1 \times 2^{4}) + (0 \times 2^{3}) + (0 \times 2^{2}) \\
       &    & + (1 \times 2^{1}) + (1 \times 2^{0})
\end{array}\]
is written: \[10010011\] in base two. More generally, the
 number written: \[a_{n}a_{n-1} \ldots a_{1}a_{0}\]
 in base two (where each of $a_{n} \ldots a_{0}$ is 0 or 1)
 is equal to:
\[(a_{n} \times 2^{n}) + (a_{n-1} \times 2^{n-1}) + \ldots
+ (a_{1} \times 2^{1}) + (a_{0} \times 2^{0})\]
(where $2^{n}$ is equal to $2$ multiplied by itself $n$ times if $n>0$, and $2^{0} = 1$).
With this convention we can therefore code
 on one byte the natural numbers between 0 (00000000) and 255 (11111111).

\subsection{Characters ({\sc ascii} code)}
\index{Character}
 The \python\ programs we write are sequences of characters. Each time
 we press a key on the keyboard, the code of the corresponding character is entered into the
 computer's memory. The characters corresponding to the digits, to the lower-case
 or upper-case letters and to common characters such as $+$, $\ast$,
 $($, $)$, \ldots\ are called
 ``printable'' characters. Some other characters, which
cannot be written, are called ``special'' characters. This is the case, for
example, of the ``control character'' and the ``escape character'',
which are obtained by pressing the keys
%touche Ctrl
\begin{picture}(22,12)(0, 2)
\put(0,0){\framebox(22,12)[lb]{}}
\put(1.5,2){\makebox(0,0)[lb]{\tiny Ctrl}}
\end{picture}\hspace{0.3ex}
and
%Esc (touche d'echappement)
\begin{picture}(14,12)(0, 2)
\put(0,0){\framebox(14,12)[lb]{}}
\put(1.5,2){\makebox(0,0)[lb]{\tiny Esc}}
\end{picture}\hspace{0.3ex}
respectively.
These characters (printable or special) are often represented on one byte (8 bits) according to
an international code called the {\sc ascii} code. The printable characters whose
{\sc ascii} code lies between 32 and
126 are given in table \ref{tab:ascii} below. For example, the {\sc ascii} code 65 of the
character {\cha $\scriptstyle '$A$\scriptstyle '$} is obtained from this table by adding the number 64 of the row
containing the character {\cha $\scriptstyle '$A$\scriptstyle '$} to the number 1 of the column
of this character.

\index{ASCII code@{\sc ascii} code}\begin{table}[htbp]
\begin{center}
\begin{tabular}{@{}|c|c|c|c|c|c|c|c|c|c|c|c|c|c|c|c|c|@{}} \cline{2-17}
\multicolumn{1}{c|}{$+$}       &  \tiny \ 0  & \tiny \ 1 & \tiny \ 2 &\tiny \ 3 &\tiny \ 4 &\tiny \ 5 &\tiny \ 6 &\tiny \ 7
        &\tiny \ 8 &\tiny \ 9 & \tiny 10 &\tiny 11&\tiny 12 &\tiny 13 &\tiny 14 &\tiny 15 \\ \hline
\tiny 32    &  \blanc & \texttt{!} & \texttt{"} & \texttt{\#} & \texttt{\$} & \texttt{\%} & \texttt{\&} & \texttt{'}
        & \texttt{(} & \texttt{)} & \texttt{*} & \texttt{+} & \texttt{,} & \texttt{-} & \texttt{.} & \texttt{/} \\ \hline
\tiny 48    & \texttt{0} & \texttt{1} & \texttt{2} & \texttt{3} & \texttt{4} & \texttt{5} & \texttt{6} & \texttt{7}
        & \texttt{8} & \texttt{9} & \texttt{:} & \texttt{;} & \texttt{<} & \texttt{=} & \texttt{>} & \texttt{?} \\ \hline
\tiny 64    & \texttt{@} & \texttt{A} & \texttt{B} & \texttt{C} & \texttt{D} & \texttt{E} & \texttt{F} & \texttt{G}
         & \texttt{H} & \texttt{I} & \texttt{J} & \texttt{K} & \texttt{L} & \texttt{M} & \texttt{N} & \texttt{O} \\ \hline
\tiny 80     & \texttt{P} & \texttt{Q} & \texttt{R} & \texttt{S} & \texttt{T} & \texttt{U} & \texttt{V} & \texttt{W}
         & \texttt{X} & \texttt{Y} & \texttt{Z} & \texttt{[} & \texttt{$\backslash$} & \texttt{]} & \texttt{\^{ } } & \texttt{\_} \\ \hline
\tiny 96     & \texttt{`} & \texttt{a} & \texttt{b} & \texttt{c} & \texttt{d} & \texttt{e} & \texttt{f} & \texttt{g}
         & \texttt{h} & \texttt{i} & \texttt{j} & \texttt{k} & \texttt{l} & \texttt{m} & \texttt{n} & \texttt{o} \\ \hline
\tiny 112  & \texttt{p} & \texttt{q} & \texttt{r} & \texttt{s} & \texttt{t} & \texttt{u} & \texttt{v} & \texttt{w}
         & \texttt{x} & \texttt{y} & \texttt{z} & \texttt{\{} & \texttt{|} & \texttt{\}} & \texttt{\~{ }} &  \multicolumn{1}{c}{}\\ \cline{1-16}
\end{tabular}
\end{center}
\caption{Characters with {\sc ascii} codes between 32 and 126\label{tab:ascii}}
\end{table}

When we say that the character {\cha $\scriptstyle '$A$\scriptstyle '$} has {\sc ascii} code 65, we mean
that this character is represented in the computer's memory by the byte
01000001, which also represents the number 65 written in base two. Notice that
the character zero has code 48. The integer 0 (coded in the machine on one byte as 00000000) therefore
has a code different from that of the character zero {\cha $\scriptstyle '$0$\scriptstyle '$} (coded in the machine as 00110000). Table
\ref{tab:asciispec} (page \pageref{tab:asciispec}) shows the {\sc ascii} codes of a few
special characters.

 \begin{table}[htbp]
\begin{center}
\begin{tabular}{@{}|c|c|c|c|c|@{}} \hline
{\sc ascii} code & 8 &  13 & 27 & 127 \\ \hline
Keyboard key \rule[-3mm]{0mm}{8mm}&
%touche d'effacement 8
\framebox[11.mm][l]{$\leftarrow$~~~}
&
\framebox[11.mm][l]{$\hookleftarrow$~~~}
&
%Esc (touche d'echappement) 27
\begin{picture}(14,12)(0, 2)
\put(0,-1){\framebox(14,12)[lb]{}}
\put(1.5,2){\makebox(0,0)[lb]{\tiny Esc}}
\end{picture}
&
%touche Del
\begin{picture}(14,12)(0, 2)
\put(0,-1){\framebox(14,12)[lb]{}}
\put(1.5,2){\makebox(0,0)[lb]{\tiny Del}}
\end{picture}
\\ \hline
\end{tabular}
\end{center}
\caption{{\sc ascii} codes of a few special characters\label{tab:asciispec}}
\end{table}

Accented letters such as {\cha $\scriptstyle '$é$\scriptstyle '$} or
{\cha $\scriptstyle '$ç$\scriptstyle '$} (used in French) do not appear in the {\sc ascii} code, which was designed
for the English language. Today they are coded according to the international code
{\sc unicode}, which can represent the characters of all the world's writing systems.
In the {\sc utf-8} format used by \python\ for the text of programs, {\sc ascii} characters are coded
on one byte as above, and the other characters on two, three or four bytes.
For example the letter {\cha $\scriptstyle '$é$\scriptstyle '$} is coded by the two bytes
11000011 10101001.

\subsection{Text}

A text (such as a \python\ program) is a sequence of characters. The code
of a text in the machine is the sequence of the codes of the characters that make it up. For example
the text {\cha "Python"} will be represented in the machine by the sequence of the {\sc ascii} codes of the characters
 {\cha $\scriptstyle '$P$\scriptstyle '$}, {\cha $\scriptstyle '$y$\scriptstyle '$}, {\cha $\scriptstyle '$t$\scriptstyle '$},
 {\cha $\scriptstyle '$h$\scriptstyle '$}, {\cha $\scriptstyle '$o$\scriptstyle '$}, {\cha $\scriptstyle '$n$\scriptstyle '$},
  that is, according to table \ref{tab:ascii}
(page \pageref{tab:ascii}),
the sequence 80 (01010000),
121 (01111001), 116 (01110100), 104 (01101000), 111 (01101111), 110 (01101110).
The code of the text {\cha "Python"} is therefore
01010000 01111001 01110100 01101000 01101111 01101110.

\subsection{Character strings}
\index{Character string}

A program can make the computer handle booleans and integers, but
also texts, all coded in binary. So as not to confuse the text of the program
with the texts handled by the program, we shall call the latter ``character
strings''.

\subsubsection{Use and notation in \protect\python:}
In \python\ we can handle texts as long as we like, which are called
``character strings''. They have type {\idt str} (short for {\idt string}). A character string is written between
straight double quotes, as for example:
\begin{center}
{\cha "This\ is\ a\ character\ string"}
\end{center}
While it is easy to pick out an identifier or an integer in the text of a
program, it is harder to recognise a text inside this text! The quotes therefore
mark off the part of the text that forms the character string from the rest of the program. But this causes
a problem if the character string contains a double quote. That is why the convention is to put
a backslash {\cha \char92} in front of the double quotes inside a character string,
as in:
\begin{center}
{\cha "The\ \char92"Python\char92"\ language"}.
\end{center}
A character string consisting of a single double quote is therefore written
{\cha "\char92""} and a string consisting of a single backslash is written
{\cha "\char92\char92"}. In \python, a character string can also be written between
single quotes (apostrophes), as in {\cha \textquotesingle{}The\ "Python"\ language\textquotesingle{}}, which avoids
having to put a backslash in front of the double quotes. \python\ has no special type
for characters (\underline{char}acters, often called {\em char} in other languages): the character
{\cha $\scriptstyle '$A$\scriptstyle '$} is simply the character string {\cha "A"}, of length 1.

\subsubsection{Concatenation of character strings:}
\index{Concatenation}
If $c_{1}$ and $c_{2}$ are character
strings then $c_{1}$\ {\idt +}\ $c_{2}$ is the character
string obtained by concatenating (that is, putting end to end) the character
strings $c_{1}$ and $c_{2}$. Several strings can be concatenated:
$c_{1}$\ {\idt +}\ $c_{2}$\ {\idt +}\ \ldots\ {\idt +}\ $c_{n}$ (so the {\idt +} sign, which
adds two numbers, puts two character strings end to end). For example:
\begin{center}
{\cha "The\ "}\ {\idt +}\ {\cha "\char92""}\ {\idt +}\ {\cha "Python"}\ {\idt +}\ {\cha "\char92""}\
{\idt +}\ {\cha "\ language"}
\end{center}
is equal to the character string that is also written:
\begin{center}
{\cha "The\ \char92"Python\char92"\ language"}
\end{center}

\subsubsection{Writing on the robot's drawing and below it:}
Whether its pen is up or down, the \robot\ can write character strings,
without moving and without changing its orientation. It can write in the drawing window or in the
message box located below the drawing window.

The command {\idt label($c$)} tells the \robot\ to write the character string $c$ to the right of the point
where it stands in its drawing window.
The command {\idt message($c$)} tells the \robot\ to write the character string $c$ in the
box provided below the window in which it moves. This erases the previous message
displayed in this
box. To give the user time to read the message, one can, if necessary, use the procedure
{\idt markPause()}, which waits until the user presses a key or clicks the mouse
button before continuing. The command {\idt clearMessage()} erases this message. The commands {\idt textFont},
{\idt textStyle} and {\idt textSize} of the {\idt LibRobot} library let you choose a font,
a style or a character size other than the standard one (on some computers, they have
no effect).

If $n$ is an integer then {\idt str($n$)} is the character
string denoting this integer in the usual decimal
notation (the {\idt LibRobot} library also offers {\idt intToString($n$)},
which does the same thing). If $b$ is a boolean then {\idt boolToString($b$)} is
{\cha "True"}
if $b$ is equal to {\mcl True}, and
{\cha "False"}
if $b$ is equal to {\mcl False} (as does {\idt str($b$)}).
For example if {\idt x} is an integer parameter of a procedure equal to 17 then the
command:
\begin{center}
{\idt label(}{\cha "The\ value\ of\ x\ is\ "}\ {\idt +\ str(x)\ +}\ {\cha "."}{\idt )}
\end{center}
writes the concatenation of the character strings
{\cha "The\ value\ of\ x\ is\ "},
{\cha "17"}
and {\cha "."}, that is:
\begin{center}
{\idt The value of x is 17.}
\end{center}
to the right of the \robot. Notice that the conversion {\idt str(x)} is essential: without it,
\python\ would refuse to add an integer to a character string and would stop the program with
the message {\idt TypeError: can only concatenate str (not "int") to str}.
If {\idt b} is a boolean parameter of a procedure equal to {\mcl True} then the
command:
\begin{center}
{\idt message(}{\cha "b\ =\ "}\ {\idt +\ boolToString(b))}
\end{center}
writes the concatenation of the character strings
{\cha "b\ =\ "}
and {\cha "True"}, that is:
\begin{center}
{\idt b = True}
\end{center}
in the message box.


\section{Music}

This section is meant for musicians who know how to read music, that is, the way
music is coded. The \robot\ can play music, in a
rather limited way it must be said, since it can only play the piano with one finger. To
code music for the \robot\ we have procedures with parameters. So all we
need is to set up the correspondence with classical music theory
\cite{solfege} to define how to write music for the \robot.

\subsection{Note values}

The procedure {\idt note(p, d)} plays
on the computer's loudspeaker a note of pitch {\idt p} and
value (duration) {\idt d}. The pitch of the note is coded by an integer {\idt p} as shown in figure
\ref{fig:codenotes} (page \pageref{fig:codenotes}).
\begin{figure}[htbp]
%\includevspacePostScript[13.4cm]{CodageNotes}{0.95}{0.95}
 \begin{flushleft}
\parindent0pt
\noindent Octave 2:\\
\begin{music}
\parindent0pt
\startextract
\NOtes\ql{q}\ql{^q}\ql{r}\ql{^r}\ql{s}\ql{t}\ql{^t}\ql{u}\ql{^u}\ql{v}\ql{^v}\ql{w}\ql{x}\enotes
\endextract
\end{music}
\vspace*{-1.21em}%
\def\sauth{\hspace*{1.19em}}%
\hspace*{17mm}{\small24\sauth25\sauth26\sauth27\sauth28\sauth29\sauth30\sauth31\sauth32\sauth33\sauth34\sauth35\sauth36}\\[1ex]
\noindent Octave 1:\\
\begin{music}
\parindent0pt
\startextract
\NOtes\ql{j}\ql{^j}\ql{k}\ql{^k}\ql{l}\ql{m}\ql{^m}\ql{^n}\ql{n}\ql{o}\ql{^o}\ql{p}\ql{q}\enotes
\endextract
\end{music}
\vspace*{-1.em}%
\def\sauth{\hspace*{1.19em}}%
\hspace*{17mm}{\small12\sauth13\sauth14\sauth15\sauth16\sauth17\sauth18\sauth19\sauth20\sauth21\sauth22\sauth23\sauth24}\\[1ex]
\noindent Octave 0:\\
\begin{music}
\parindent0pt
\startextract
\NOtes\qu{c}\qu{^c}\qu{d}\qu{^d}\qu{e}\qu{f}\qu{^f}\qu{g}\qu{^g}\qu{h}\qu{^h}\ql{i}\ql{j}\enotes
\endextract
\end{music}
\vspace*{-0.75em}%
\def\sauth{\hspace*{1.67em}}%
\hspace*{18mm}{\small0\sauth1\sauth2\sauth3\sauth4\sauth5\sauth 6\sauth7\sauth8\sauth9\sauth1\rlap{0}\sauth1\rlap{1}\sauth1\rlap{2}}\\[1ex]
\noindent Octave -1:\\
\begin{music}
\parindent0pt
\setclef1\bass
\startextract
\NOtes\qu{J}\qu{^J}\qu{K}\qu{^K}\qu{L}\qu{M}\qu{^M}\qu{N}\qu{^N}\qu{a}\qu{^a}\ql{b}\ql{c}\enotes
\endextract
\end{music}
\vspace*{-1.2em}%
\def\sauth{\hspace*{1.445em}}%
\hspace*{16mm}{\small-1\rlap{2}\sauth-1\rlap{1}\sauth-1\rlap{0}\sauth -9\sauth -8\sauth -7\sauth -6\sauth -5\sauth -4\sauth -3\sauth -2\sauth -1\sauth 0}\\[1ex]
\noindent Octave -2:\\
\begin{music}
\parindent0pt
\setclef1\bass
\startextract
\NOtes\qu{C}\qu{^C}\qu{D}\qu{^D}\qu{E}\qu{F}\qu{^F}\qu{G}\qu{^G}\qu{H}\qu{^H}\qu{I}\qu{J}\enotes
\endextract
\end{music}
\def\lignes#1#2{\begin{tabular}[t]{@{}c@{}}#1\\#2\end{tabular}}
\def\sauth{\hspace*{0.86em}}%
\hspace*{16mm}{\small\lignes{-24}{C}\sauth -23\sauth \lignes{-22}{D}\sauth-21\sauth \lignes{-20}{E}\sauth\lignes{-19}{F}\sauth -18\sauth \lignes{-17}{G}\sauth -16\sauth\lignes{-15}{A}\sauth -14\sauth\lignes{-13}{B}\sauth\lignes{-12}{C}}
\end{flushleft}%
 \caption{Coding the pitch of musical notes by integers}\label{fig:codenotes}
\end{figure}
The value of the note gives its duration, which is coded by a constant {\idt d} that can be {\idt whole} (whole note),
{\idt dotted\_whole} (dotted whole note),
{\idt half} (half note),
{\idt dotted\_half} (dotted half note),
{\idt half\_triplet} (half-note triplet),
{\idt quarter} (quarter note),
{\idt dot\-ted\_\-quar\-ter} (dotted quarter note),
{\idt quar\-ter\_\-trip\-let} (quarter-note triplet),
 {\idt eighth} (eighth note),
 {\idt dotted\_eighth} (dotted eighth note),
{\idt eighth\_\-trip\-let} (eighth-note triplet),
{\idt sixteenth} (sixteenth note),
{\idt dot\-ted\_\-six\-teenth} (dotted sixteenth note) or
{\idt thir\-ty\_\-sec\-ond} (thirty-second note). These constants, declared in the {\idt LibRobot} library,
are numbers written with a decimal point (as we shall explain in chapter 19), which
give the duration of the note in beats: a {\idt quarter} is worth {\idt 1.0},
an {\idt eighth} is worth {\idt 0.5}, a {\idt whole} is worth {\idt 4.0}, etc.
The relative duration of a {\idt dot\-ted\_\-whole} is
one and a half times that of a
 {\idt whole}. The relative duration of a {\idt whole} is twice that of a
 {\idt half}, three times that of a {\idt half\_\-trip\-let}, \ldots\
Other note values can be defined, for example
{\idt sixty\_fourth = 0.0625} (sixty-fourth note).

 \subsection{Tempo}

 The tempo,
 that is, the degree of speed or slowness with which the piece of music must be played,
 determines the absolute duration of the different note values. The tempo is set
 by the procedure {\idt metronome($n$)}, which gives the number $n$ of beats of the
 metronome per minute, a quarter note being played on each beat. For example
 {\idt metronome(60)} means that 60 quarter notes must be played per minute. Tempos
 are usually given by Italian terms (adagio, allegro, prestissimo, \ldots) which
 are placed at the beginning of a piece, above the staff. They are defined
 in the {\idt LibRobot} library by the following constant declarations:

\input python/Tempos.tex

\subsection{Notes}

\noindent The numbering of the notes in figure \ref{fig:codenotes} (page \pageref{fig:codenotes}) is
not easy to remember. That is why we shall code the notes by
an integer expression over constants that are easy to remember.

A note is defined by its name $n$ (namely C, D, E, F, G, A, B), its possible
accidental $a$ (that is, a sign that changes the sound of the note it is attached to,
which can be the flat $\flat$ (which lowers the note by a semitone),
the sharp $\sharp$ (which raises the note by a semitone) or
the natural $\natural$ (which cancels the effect of the sharp or the flat)) and
its octave $o$ (which tells to which series of notes with the same name this note belongs).
We code the octaves from low to high by the integers -3, -2, \ldots, 3, 4,
octave 0 starting with the C (middle C) at the centre of the piano keyboard. The
following constants are declared in the {\idt LibRobot} library:

\input python/Notes.tex

\noindent The number of the note as defined in figure \ref{fig:codenotes}
(page \pageref{fig:codenotes}) is given by the value of the integer expression
$o \ast {\idt octave} + n + a$. If the note has no accidental, its code is simply
$o \ast {\idt octave} + n$. For example the C in the middle of the piano keyboard corresponds to
$o = 0$, $n = ${\idt C}, that is $0 \ast ${\idt octave}$ + 0 = 0$. The D sharp one octave higher
has code $1 \ast ${\idt octave}$ + ${\idt D}$ + ${\idt sharp}$
= 1 \ast 12 + 2 + 1 = 15$. The neighbouring D flat
has code $1 \ast ${\idt octave}$ + ${\idt D}$ + ${\idt flat}$
= 1 \ast 12 + 2 + (-1) = 13$. When the pitch of a note is given by an expression,
we simply write this expression as the first parameter of {\idt note}, as in {\idt note(D + octave + sharp, quarter)}.

\subsection{Rests}

The procedure {\idt rest($s$)} interrupts the sound for a relative duration set by
the rest value $s$, which can be one of the constants
  {\idt whole\_rest},
  {\idt dotted\_half\_rest},
  {\idt half\_rest},
  {\idt quarter\_rest},
  {\idt dot\-ted\_\-quar\-ter\_\-rest}\footnote{\ It is not customary to dot the quarter rest.},
  {\idt eighth\_\-rest},
  {\idt dotted\_eighth\_rest},
  {\idt sixteenth\_rest}
or  {\idt thirty\_second\_rest}.

\subsection{Dynamics}

Dynamics are the different degrees of loudness through which one or more sounds, a passage
or a whole piece can go. The procedure {\idt dynamics($n$)} sets the dynamics with which
the rest of the piece must be played. The possible dynamics terms $n$ are, from the softest to the loudest:
{\idt pia\-nis\-sis\-si\-mo},
{\idt pia\-nis\-si\-mo},
{\idt piano},
{\idt mezzo\_piano} (or {\idt mezza\_voce}),
{\idt mezzo\_forte},
{\idt forte},
{\idt fortissimo} or
{\idt for\-tis\-sis\-si\-mo}.

\begin{exemple}[Au clair de la lune]{\rm
``Au clair de la lune'' (``By the light of the moon''), said to have been composed by the musician Jean-Baptiste Lully,
is probably the most famous of all French children's songs:

%\includevspacePostScript[4.7cm]{AuClairDeLaLune}{0.95}{0.95}
%%&frLaTeX
%\documentclass{article}
%\usepackage{musixtex}
%\begin{document}
\begin{flushleft}
\parindent0pt
\begin{music}
\generalmeter{\meterC}
\generalsignature{1}
\nostartrule
\parindent0pt%
\nobarnumbers
\startpiece
\notesp\qu{g}\qu{g}\qu{g}\qu{h}\enotes\bar
\notesp\hl{i}\hu{h}\enotes\bar
\notesp\qu{g}\ql{i}\qu{h}\qu{h}\enotes\Setvolta1\bar
\notesp\zw{g}\enotes\setendvoltabox\Setvolta2\rightrepeat%
\notesp\zw{g}\enotes\setendvoltabox\alaligne
\NOtesp\qu{h}\qu{h}\qu{h}\qu{h}\enotes\bar
\NOtesp\hu{e}\hu{e}\enotes\bar
\NOtesp\qu{h}\ql{g}\qu{f}\qu{e}\enotes\bar
\NOtesp\zw{d}\enotes\alaligne
\NOtesp\qu{g}\qu{g}\qu{g}\qu{h}\enotes\bar
\NOtesp\hl{i}\hu{h}\enotes\bar
\NOtesp\qu{g}\ql{i}\qu{h}\qu{h}\enotes\bar
\NOtesp\zw{g}\enotes
\Endpiece
\end{music}
\end{flushleft}%
%\end{document}%

\noindent The program below makes the computer play it. To
save space, we have written two notes per line, with the corresponding syllables
of the (French) lyrics in comments (those of the first and second lines of verse are separated by a
slash {\cmt /}):

 \input python/AuClairDeLaLune.tex
\noindent $\Box$ }\end{exemple}


\newpage
\needspace{6\baselineskip}\noindent \rule{137mm}{0.1mm}\par\nopagebreak %---------------------------------------------
\begin{exercices}{\ \\
\vspace{-\baselineskip}

\sloppy
\tiret  Write procedures {\idt startOfTransmission}, {\idt a}, \ldots, {\idt z}, {\idt c0},
\ldots, {\idt c9} (digits), {\idt fullStop}, {\idt comma},
 {\idt semicolon}, {\idt colon}, {\idt questionMark},
 {\idt exclamationMark}, {\idt hyphen}, {\idt apostrophe},
  {\idt slash},
 {\idt error}, {\idt endOfTransmission}, {\idt endOfWord} and
 {\idt newLine} (go to the next line), then use them to code
 messages in Morse, each sign of the plain message being coded by calling one
 of these procedures. For example ``{\em coded message}'' will be translated into the sequence of instructions
`` {\ttfamily\idt c(); o(); d(); e(); d(); endOfWord(); newLine();
 m(); e(); s(); s(); a(); g(); e(); endOfWord()}''. You can add sound by using the following procedures
 of the \robot:
 \begin{list}{$\scriptscriptstyle\bullet$}{
}
   \item {\idt beep()} produces a short sound;
   \item {\idt longBeep()} produces a long sound;
   \item {\idt delay({\rm\em t})} waits for {\em t} seconds ({\em t} can be a number written with a decimal point, as in {\idt delay(0.5)}).
 \end{list}

 \tiret Before the invention of electricity, France used the optical (semaphore) telegraph invented in
 1793 by Citizen Claude Chappe. The telegraph consists of three articulated wooden arms
 at the top of a mast set up on a tower on a hilltop or on a church steeple.
 As soon as he sees through his telescope a sign of activity at the nearest station, the {\em telegraph operator}
 hurries to work the arms of his own machine in turn to pass on the message (without understanding it)
  to the next station. Step by step the message reaches its destination. It is in a
  way an improved version of the smoke signals of the American Indians! After a first
  trial on 12 July 1793, the first line was opened between Paris and Lille on 17 August 1794
  \cite{bertho}. The signals of Chappe's optical telegraph could be decoded
  according to several codes. The alphabetic code is given in figure \ref{fig:chappe} (page
  \pageref{fig:chappe}). Write a program that displays the signals successively, one by one,
  and sends the message ``Paris est tranquille et les bons citoyens sont contents''
  (``Paris is calm and the good citizens are content''),
  which Bonaparte had broadcast immediately after the coup d'état of 18 Brumaire (9--10 November
  1799).

  To leave the signal visible for some time, use the procedure
  {\idt delay({\rm\em t})}, which waits for {\rm\em t} seconds. To clear the screen
  between signals, use the procedure {\idt bg(white)}, which paints the background of the screen
  white (or another colour, depending on the parameter) and puts the \robot\ back in its initial state
  (visible at the centre of the screen, facing north, black pen of width 1 down, with a movement
  grid of (12~$\times$~12)).

  \tiret Write a \python\ program that translates the letters of the alphabet into the phonetic alphabet
  according to the international code of table \ref{tab:alphone} (page \pageref{tab:alphone}).

   \tiret Write a \python\ program that encrypts and decrypts messages using a simple
   secret code, obtained for example by a permutation of the letters of the alphabet.

\noindent $\Box$}\end{exercices}
\begin{exercices}{\ \\
\vspace{-\baselineskip}

\tiret What is the 16-bit binary representation of the number written 127
in base 10?

\tiret What is the decimal notation of the positive integer represented
in binary on 16 bits by 0011010101010101?

\tiret What text is represented in {\sc ascii} code by the following bytes:
01000100 01100101 01110011 01100011 01100001
01110010 0111\-0100 0110\-0101 0111\-0011?

\noindent $\Box$}\end{exercices}
\begin{figure}[htbp]
 \includevspacePostScript[9.6cm]{TelegrapheChappe}{0.95}{0.95}
 \caption{Alphabet of Chappe's optical telegraph}\label{fig:chappe}
\end{figure}
\begin{table}[htbp]
\begin{center}
\begin{tabular}{@{}|l|l||l|l||l|l||l|l||l|l|@{}} \hline
A & Alpha   & B  & Bravo & C & Charlie & D & Delta   & E  & Echo \\\hline
F & Foxtrot  & G & Golf   & H  & Hotel & I & India & J & Juliet \\\hline
K  & Kilo & L & Lima & M & Mike   & N  & November & 0 & Oscar \\\hline
P & Papa   & Q  & Quebec & R & Romeo  & S & Sierra   & T  & Tango \\\hline
U & Uniform & V & Victor   & W  & Whiskey & X & Xray & Y & Yankee   \\\hline
Z  & Zulu &   \multicolumn{8}{c}{}   \\\cline{1-2}
\end{tabular}
\end{center}
\caption{International phonetic alphabet}\label{tab:alphone}
\end{table}
\begin{exercice}{Songs are an inexhaustible source of exercises
 \cite{boutet}, \cite{sabatier1}, \cite{sabatier2}, \cite{sabatier3}.
 For example, write
\python\ programs to play {\em Green\-sleeves\/} (figure \ref{fig:Greensleeves}, page \pageref{fig:Greensleeves})
and {\em O Tannenbaum\/} (figure \ref{fig:Tannenbaum}, page \pageref{fig:Tannenbaum}).
You can also scroll the score line by line on the screen
while it is being played by the computer, drawing the staff
and the notes as in figure \ref{fig:RepNotes} (page \pageref{fig:RepNotes}).
\noindent $\Box$}\end{exercice}
\begin{figure}[htbp]
 %\includevspacePostScript[12.4cm]{Greensleeves}{0.95}{0.95}
 %%&frLaTeX
%\documentclass{article}
%\usepackage{musixtex}
%\begin{document}
\begin{flushleft}
\parindent0pt
\begin{music}
\generalmeter{\meterfrac34}
\nostartrule
\parindent0pt%
\nobarnumbers
\startpiece
%
\NOtes
\Uptext{\metron{\qu}{144} Allegro}%
\qp\qp\qu{h}\enotes\doublebar
\NOtes\hl{j}\ql{k}\enotes\bar
\NOtes\qlp{l}\cl{^m}\ql{l}\enotes\bar
\NOtes\hl{k}\ql{i}\enotes\bar
\NOtes\qup{g}\cu{h}\ql{i}\enotes\bar
%
\NOtes\hl{j}\qu{h}\enotes\bar
\NOtes\qup{h}\cu{^g}\qu{h}\enotes\bar
\NOtes\hl{i}\qu{^g}\enotes\bar
\NOtes\hu{e}\qu{h}\enotes\bar
\NOtes\hl{j}\ql{k}\enotes\bar
\NOtes\qlp{l}\cl{^m}\ql{l}\enotes\bar
%
\NOtes\hl{k}\ql{i}\enotes\bar
\NOtes\qup{g}\cu{h}\qu{i}\enotes\bar
\NOtes\qlp{j}\cl{i}\qu{h}\enotes\bar
\NOtes\qup{^g}\cu{^f}\qu{g}\enotes\bar
\NOtes\hup{h}\enotes\bar
%
\NOtes\qu{h}\ibl013\qb0{h}\qb0{i}\qb0{j}\tbl0\qb0{l}\enotes\leftrepeat
\NOtes\hlp{n}\enotes\bar
\NOtes\qlp{n}\cl{^m}\ql{l}\enotes\bar
\NOtes\hl{k}\ql{i}\enotes\bar
\NOtes\qup{g}\cu{h}\ql{i}\enotes\bar
\NOtes\hl{j}\qu{h}\enotes\bar
%
\NOtes\qup{h}\cu{^g}\qu{h}\enotes\bar
\NOtes\hl{i}\qu{^g}\enotes\bar
\NOtes\hup{e}\enotes\bar
\NOtes\hup{n}\enotes\bar
\NOtes\qlp{n}\cl{^m}\ql{l}\enotes\bar
\NOtes\hl{k}\ql{i}\enotes\bar
%
\NOtes\qup{g}\cu{h}\ql{i}\enotes\bar
\NOtes\qlp{j}\cl{i}\qu{h}\enotes\bar
\NOtes\qup{^g}\cu{^f}\qu{g}\enotes\bar
\NOtes\hup{h}\enotes\Setvolta1\bar
\NOtes\qu{h}\ibl013\qb0{h}\qb0{i}\qb0{j}\tbl0\qb0{l}\enotes\setendvoltabox\Setvolta2\rightrepeat%
\NOtes\hu{h}\qp\enotes\setendvoltabox
%
\Endpiece
\end{music}
\end{flushleft}%
%\end{document}%
\endinput
 \caption{Greensleeves (Irish song)}\label{fig:Greensleeves}
\end{figure}
\begin{figure}[htbp]
 %\includevspacePostScript[2.8cm]{OTannenbaum}{0.95}{0.95}
 %%&frLaTeX
%\documentclass{article}
%\usepackage{musixtex}
%\begin{document}
\begin{flushleft}
\parindent0pt
\begin{music}
\generalmeter{\meterfrac34}
\nostartrule
\parindent0pt%
\nobarnumbers
\startpiece
%
\notes\qp\qp\qu{c}\enotes\bar
\notes\cup{f}\ccu{f}\qu{f}\qu{g}\enotes\bar
\notes\cup{h}\ccu{h}\qu{h}\qu{f}\enotes\bar
\notes\cu{g}\cu{h}\ql{i}\qu{e}\enotes\bar
\notes\qu{g}\qu{f}\enotes\setrightrepeat\alaligne
%
\notes\qp\qp\ds\cl{j}\enotes\bar
\notes\cl{j}\cu{h}\qlp{k}\cl{j}\enotes\bar
\notes\cl{j}\cl{i}\qlp{i}\cl{i}\enotes\bar
\notes\cl{i}\cl{g}\qlp{j}\cl{i}\enotes\bar
\notes\cl{i}\cu{h}\qu{h}\enotes\setrightrepeat
%
\endpiece
\end{music}
\end{flushleft}%
%\end{document}%
\endinput
 \caption{O Tannenbaum (``O Christmas Tree'', German song)}\label{fig:Tannenbaum}
\end{figure}
\begin{figure}[htbp]
 \includevspacePostScript[2.3cm]{EcrireEtJouerMus}{0.95}{0.95}
 \caption{Staff and notes drawn on the computer screen}\label{fig:RepNotes}
\end{figure}
\begin{corrige}{
We give one program, the others being in the folder {\idt book/\allowbreak 18\_coding}:
\input python/MorseCode1.tex

The same folder contains various possible improvements of this program, in particular
automatic line wrapping ({\idt MorseCode2}).
The procedures {\idt a}, \ldots, {\idt z} and {\idt c0}, \ldots, {\idt c9} have
lower-case names because the custom in \python\ is to begin procedure names with a lower-case letter,
names beginning with a capital letter usually being reserved for something else.
\noindent $\Box$}\end{corrige}
\begin{corriges}{\ \\
\vspace{-\baselineskip}

\tiret The 16-bit binary representation of 127 is
00\-00\-00\-00\-01\-11\-11\-11.

\tiret The number represented by the half-word 0011010101010101 is written
13653 in base 10.

\tiret The sequence of bytes 01000100 01100101 01110011 01\-10\-00\-11 01\-10\-00\-01
01\-11\-00\-10 01110100 01100101 01110011,
that is, in decimal notation, 68 101 115 99 97 114 116 101 115,
represents the character string
{\cha $\scriptstyle '$Descartes\idt $\scriptstyle '$}.

\noindent $\Box$}\end{corriges}
\begin{corrige}{
We give one program, many other songs being programmed in the folder
{\idt book/\allowbreak 18\_coding} and in its sub-folder {\idt songs}:
\input python/OTannenbaum.tex
\noindent $\Box$}\end{corrige}

%%%%%%%%%%%%%%%%%%%%%%%%%%%%%%%%%%%%%%%%%%%%%%%%%%%%%%%%%%%%%%%%%%%%%%%%%%%
\noindent \framebox[\textwidth]{
\begin{minipage}{\myparboxwidth}
\setlength{\parindent}{2em}
\small \sf
\vspace{\abovedisplayskip}
\section*{Addenda}

In addition, we explain two representations of signed integers, the second
being the more widely used.

\paragraph*{Sign-and-magnitude representation of signed integers:\ }
\index{Binary number system}
To code a signed integer (positive or negative) on $n$ bits, one can reserve one bit
for the sign and code the absolute value of the integer on the remaining $n - 1$ bits.
Let us adopt the convention
of putting the sign first, choosing 0 for the sign $+$ and 1 for the sign $-$.

\vspace{\abovedisplayskip}
\end{minipage}
}
%%%%%%%%%%%%%%%%%%%%%%%%%%%%%%%%%%%%%%%%%%%%%%%%%%%%%%%%%%%%%%%%%%%%%%%%%%%%

%%%%%%%%%%%%%%%%%%%%%%%%%%%%%%%%%%%%%%%%%%%%%%%%%%%%%%%%%%%%%%%%%%%%%%%%%%%
\noindent \framebox[\textwidth]{
\begin{minipage}{\myparboxwidth}
\setlength{\parindent}{2em}
\small \sf
\vspace{\abovedisplayskip}
The code of
$-10$ on one byte will be 10001010, that is, 1 for the sign followed by 0001010 for the absolute value $10$
of $-10$. With
this convention one can code on one byte the integers between $-127$ (11111111) and $127$
(01111111). The integer zero is represented by 00000000 ($+0$) or by 10000000 ($-0$).

\paragraph*{Two's complement representation of signed integers:\ }
In the two's complement representation on $n$ bits, the positive numbers
strictly less than $2^{n-1}$ are
represented as before in base 2 on $n-1$ bits, the first (sign) bit being
zero. A negative number $-x$
such that $x \leq 2^{n-1}$ is represented by a first sign bit equal to 1 followed by the
representation of the positive number $p < 2^{n-1}$ such that $-x = -2^{n-1} + p$,
coded
on $n - 1$ bits. For example, to code $-1$ on
one byte, we write $-1 = -128 + 127$ and code 127 on 7 bits, which gives 1111111, which we
precede by 1 to get 11111111. With this convention we can also code
$-128 = -128 + 0$, which gives the code 10000000. This method makes it possible to
code on one byte all the signed integers between $-128$ and $+127$. \index{Integer}On 16 bits, one can
code the integers between $-32768$ and $32767$, which were the integers
of the first microcomputers. On 32 bits one can code the integers between
$-2\,147\,483\,648$ and $2\,147\,483\,647$. This representation leads to very simple algorithms
for addition and subtraction. Moreover the algorithm for coding a negative number
$-x$ is also simple. We code the positive number $x$ on $n$ bits, then invert the bits,
0 into 1 and 1 into 0, then add 1. For example to code $-127$, we code 127 as 01111111, then
invert all the bits, which gives 10000000, and add 1 to get 10000001.
The integers of type {\idt int}
in \python\ are not limited to 32 or 64 bits: \python\ uses as many bits as
needed (within the limits of the computer's memory), so that {\idt 2 ** 100}
(2 multiplied by itself 100 times) is exactly
$1\,267\,650\,600\,228\,229\,401\,496\,703\,205\,376$.
\vspace{\abovedisplayskip}
\end{minipage}
}
%%%%%%%%%%%%%%%%%%%%%%%%%%%%%%%%%%%%%%%%%%%%%%%%%%%%%%%%%%%%%%%%%%%%%%%%%%%%
\newpage                 % Si on termine le chapitre par une page blanche
\thispagestyle{empty}

